\section{Selector Human-Evaluation Protocol}

\subsection{Primary question}

When choosing among multiple generated videos for the same prompt, does a
TRIBE/BMD memorability selector choose clips that humans judge as more memorable
than clips chosen by non-brain baselines?

\subsection{Experimental unit}

The experimental unit is the prompt-conditioned candidate set. Each set contains
multiple videos generated from the same prompt or image seed. Selector policies
choose one candidate per set. Human raters compare videos only within the same
set, reducing confounding from prompt-level differences.

\subsection{Selector policies}

The minimum policy set is:
\begin{enumerate}
  \item random candidate selection;
  \item CLIP or text-video prompt-preservation winner;
  \item V-JEPA memorability winner;
  \item TRIBE/BMD memorability winner;
  \item TRIBE/BMD memorability winner with preservation or quality gate.
\end{enumerate}
Optional policies include aesthetic/video-quality winners, LoRA candidates, and
ensemble policies that combine memorability, quality, and prompt preservation.

\subsection{Human task}

The first study should use a fast pairwise preference task:

\begin{quote}
Which video do you think would be more memorable if you saw many such clips?
\end{quote}

This measures predicted memorability rather than delayed recognition. If this
pilot produces a clear effect, a follow-up delayed-recognition study should
present target and filler clips, wait for a delay, and then test recognition
against repeated and lure clips.

\subsection{Primary endpoint}

The primary endpoint is the prompt-clustered win rate of TRIBE+gate selected
clips against the strongest non-brain baseline. The primary hypothesis is:
\[
P(\text{human chooses TRIBE+gate over baseline}) > 0.5.
\]

\subsection{Statistical analysis}

The analysis should report:
\begin{itemize}
  \item prompt-clustered bootstrap confidence intervals;
  \item per-prompt win-rate distributions;
  \item mixed-effects logistic regression with selector policy as a fixed effect
        and random intercepts for prompt and, if identifiable, rater;
  \item multiple-comparison correction when comparing many selector policies.
\end{itemize}

\subsection{Anti-circularity rule}

No claim of human improvement can rely on TRIBE-scored gains alone. TRIBE may
train or select candidates, but the final endpoint must be independent human
judgment or delayed recognition.

\section{IRB-Relevant Study Summary}

The proposed first study is a minimal-risk online behavioral experiment. Adult
participants recruited through Prolific will view short generated videos and
make pairwise judgments about predicted memorability. The study does not collect
clinical data, biometric data, sensitive personal histories, or deception-based
responses. Responses will be analyzed in aggregate to compare selector policies.

The main foreseeable risks are boredom, fatigue, mild visual discomfort, and the
ordinary privacy risks associated with online survey data. Participants should be
warned that short videos may contain motion or rapid visual change, and they
should be allowed to stop at any time. The study should avoid graphic, sexual,
political persuasion, medical, or otherwise sensitive stimuli unless separately
reviewed.

Data should be stored without names, emails, IP addresses, or other direct
identifiers where possible. Prolific IDs, if collected for payment or duplicate
prevention, should be stored separately from response data or replaced with
pseudonymous participant IDs before analysis.
